\section{How High Can the Score Go? Limits by Horizon}\label{sec:horizons}

We now combine the model with the estimates to bound the score under full automation and to trace a central path to 2028, 2031 and 2036. Throughout we use the free-cash-flow-after-SBC basis and $X=1$, the Rule of 40 itself. The data do not identify the spend elasticity (Section~\ref{sec:empirics}), so the central path uses $\alpha=\CenAlpha$ and the sensitivities its range $[0.25,0.45]$. Whether the limit is a Rule of 100 or a Rule of 1000 has a precise answer in the model: it depends on the margin intercept, market growth, capital intensity, demand pull and the spend elasticity, and automation moves only some of them. How fast a firm approaches the limit depends on its capacity to execute risky transitions (Section~\ref{sec:model-attainable}). Every number in this section is conditional on the scenario values it states. The replication package grounds each parameter in an assumptions register with four classes (sourced, grounded by precedent, calibrated from our own data, scenario values), and Table~\ref{tab:param-sources} of the Supplementary Appendix lists the source behind each.

\subsection{How AI enters the parameters}
\autoref{prop:ai} maps the effects of AI onto the model's parameters; the companion paper \citep{murthy2026transformation} derives the mapping from the tasks a firm automates. Automation of general, administrative and maintenance work lowers the non-discretionary cost $a$ and raises the margin intercept $m_0=\gamma-a$, but it adds the recurring costs of running AI (\autoref{def:fullcost}). Inference costs enter gross margin $\gamma$ through the compute share $c_c$. A survey of about 300 executives at software companies building AI products puts the gross margin on those products near 45\% in 2025 and projects 53\% in 2026 and 59\% in 2027 \citep{iconiq2026state}. That is the cost structure of AI products sold to customers, including talent and hosting, not the cost of automating a firm's own work; it is below the median gross margin of \MedGrossMargin\% in our core sample in 2023--25. The price of reaching a fixed level of benchmark performance has fallen by between 9 and 900 times a year depending on the threshold \citep{epoch2025inference}, but competition may pass much of that decline to customers, so its effect on $\gamma$ depends on pricing power (\autoref{lem:prices}). Whether the compute share declines depends on how fast the volume of inference per task grows against the fall in its price, for which there is no market-level estimate: the preliminary elasticity of about $-1.1$ across providers of the same open model \citep{demirer2025emerging} measures substitution between providers, and its authors report that aggregate elasticities are substantially smaller. AI can raise demand pull $\lambda h$ for firms whose products substitute for labor and market growth $G$ where spending on wages converts into spending on software, and firms that own computing capacity raise $k$. At given prices, automation raises the saturated score by $1-\tau$ points for every point of cost it removes net of the costs it adds, up to the ceiling of \autoref{cor:ceiling}; under free entry with cheap replication and close substitutes it does not, because the margin it creates is competed away (\autoref{prop:entry}).

\subsection{The ceiling under full automation}\label{sec:horizons-ceiling}
\autoref{cor:ceiling} bounds the saturated score of a firm at the full-automation cost structure: one that has automated everything it can, spends nothing on growth and holds a saturated share. It is a bound at that cost structure, not on observed scores: a firm with leaner costs than the scenario, in particular a lower compute share, can score above it. Its parameters are brackets, not estimates. Table~\ref{tab:ceiling} evaluates it at 2025 prices, and the compute share by horizon follows below.

\emph{Accountability floor.} The floor $a_{\min}$ collects the costs of being accountable that automation does not remove. They are largely fixed in dollars, so as a share of revenue they shrink with scale (\autoref{rem:assumptions}(a)). We use a band from \AminLo\% to \AminHi\% of revenue and the values \CeilAminPctFav\%, \CeilAminPctCen\% and \CeilAminPctCons\% in the favorable, central and conservative cases. The upper end was set at the fifth-percentile share of general and administrative expense of firms with \$1 to 10 billion of revenue in the wider universe of software and internet registrants; in our prepackaged-software sample that percentile is \AminGaCoreFive\%, so even the conservative case assumes that a firm of that size spends less on general and administrative expense than 95\% of the sample does today. The central value is close to the tenth-percentile share of firms with more than \$10 billion of revenue in our sample (\AminGaCoreTen\%), so it assumes that none of that expense is automated. \autoref{prop:amin} grounds the lower end from below: internal control over financial reporting is a pass-fail standard, so the floor is the cost of the cheapest control design that leaves no reasonable possibility of an undetected material misstatement, not a trade of control cost against expected loss. In the illustrative example of \autoref{rem:amin-values}, for one control objective of a firm with \$1 billion of revenue, the floor is about \AminFloor{} to \AminFloorOne\% of revenue at the two thresholds (\AminFloorLo{} to \AminFloorHi\% over the range of audit fees): direct audit and Section 404 costs of about \AminDirect\%, from a survey of SOX programs and the average audit fee of accelerated filers \citep{kpmg2025sox,auditanalytics2023fees}, a low-side choice because a public firm of that size is normally a large accelerated filer, plus about \AminBest\% for the cheapest admissible design at a remote threshold of \AminEpsCen\% a year (\AminBestOne\% at \AminEpsSens\%), a scenario value. In that example no design built only from AI layers on one provider passes, however many layers it stacks, because the provider's common-mode probability exceeds the \AminPbarFive{} per occasion at which such a design could pass at \AminEpsCen\%; better IT general controls that lower it help, more layers do not. Each further AI application in the scope of internal control adds its own IT general controls, so the floor reaches the band's lower end at about \ItgcApps{} such applications per \$1 billion of revenue (\ItgcAppsRef{} on a survey-based line); the lower end is thus conditional on how many in-scope AI applications a firm runs. The ceiling uses the band, not this floor; with the floor in place of $a_{\min}$ the central threshold would be \GstarAminFloor\% instead of \GstarCenOne\%.

\emph{Taxes, fees and compute.} The cash tax rate $\tau$ of 15 to 21\% is a scenario range up to the 21\% federal rate of 26 U.S.C.\ 11(b), and excludes state taxes. Its lower end is not set by the 15\% corporate alternative minimum tax, which applies to corporations whose adjusted financial statement income averages more than \$1 billion over three years, or \$100 million for a US member of a foreign-parented group whose group income exceeds \$1 billion \citep{usc2026sec59}; average effective rates of profitable large corporations were lower still, 9 to 16\% in 2014--18 \citep{gao2022effective}. Fees $f$ of \FeeLo{} to \FeeHi\% of revenue are card-processing rates applied to all revenue, so they overstate the fees of firms that invoice large customers. The compute share $c_c$ of a fully automated firm has not been measured, so it is a scenario parameter. At 2025 prices we use \CcCen\% of revenue in the central case and \CcFav\% and \CcCons\% in the favorable and conservative cases: labor costs of \CcLaborShare\% of revenue, including stock-based compensation, replaced by compute at a tenth of their cost in the central case, which is above both the cost of today's agent runs and the field evidence of about three cents of AI spending per dollar of displaced freelance spending \citep{stevens2026payrolls}. The scenario range runs from \CcLo\%, labor at \CcLaborShare\% of revenue replaced at the cost per completed task of today's agent runs, failed attempts included, about \KappaLo\% of the labor's cost \citep{kwa2025measuring,metr2026evalanalysis}, to \CcHi\%, the value test, under which compute can cost at most the labor it replaces. We keep a high-compute case of \CcHigh\%, roughly the cost of revenue implied by the gross margins projected for AI products in 2027 \citep{iconiq2026state}; it is within the value test, so it is a feasible, if compute-heavy, full-automation cost structure. Section~\ref{supp:model-compute} of the Supplementary Appendix gives the evidence behind these values. For comparison, the median core software firm's entire cost of revenue was \MedCostRevenue\% of revenue in 2023--25.

\emph{Full-cost floors.} Table~\ref{tab:fullcost} in Section~\ref{app:theory} of the Supplementary Appendix evaluates \autoref{lem:fullcost}. Errors, misbehavior and common-mode losses are charged on the same base as compute, the \ErrBase\% of revenue in automated labor. Security starts from today's security spending of 0.69\% of revenue, reported in a survey of chief information security officers, and scales it to full automation; maintenance applies sourced upkeep rates and observed model-retirement intervals to a scenario stock of builds; and failure losses combine measured omission and hallucination rates, with no review at full automation, with measured recall and provider-incident rates and scenario loss severities. Together the floors come to $c_{\mathrm{full}}=\CfullFav$, \CfullCen{} and \CfullCons{} points in the favorable, central and conservative cases, and \CfullHigh{} points in the high-compute case.

\begin{table}[t]
\centering
\caption{Full-automation ceiling $R_{\max}$ at full cost and 2025 prices for the Rule of 40 (equation \eqref{eq:Rmax}, $\psi=0$, $\delta_K=0$), in percentage points. The compute share is a scenario value: no source measures it for a fully automated firm}\label{tab:ceiling}
\footnotesize
\setlength{\tabcolsep}{2.5pt}
\begin{tabular}{@{}lccccccccccccc@{}}
\toprule
& & & & & & & \multicolumn{6}{c}{$(G,\ k)$} & \\
\cmidrule(lr){8-13}
Case & $c_c$ & $f$ & $a_{\min}$ & $c_{\mathrm{full}}$ & $\tau$ & Base & $(0.05,0)$ & $(0.15,0)$ & $(0.30,0)$ & $(0.15,0.5)$ & $(0.05,1.5)$ & $(0.30,1.5)$ & $G^*$\\
\midrule
Favorable & \CeilCcFav & \CeilFFav & \CeilAminFav & \CfullFav & \CeilTauFav & \MarginFav & \CeilTabFavA & \CeilTabFavB & \CeilTabFavC & \CeilTabFavD & \CeilTabFavE & \CeilTabFavF & \GstarFavOne\\
Central & \CeilCcCen & \CeilFCen & \CeilAminCen & \CfullCen & \CeilTauCen & \MarginCen & \CeilTabCenA & \CeilTabCenB & \CeilTabCenC & \CeilTabCenD & \CeilTabCenE & \CeilTabCenF & \GstarCenOne\\
Conservative & \CeilCcCons & \CeilFCons & \CeilAminCons & \CfullCons & \CeilTauCons & \MarginCons & \CeilTabConsA & \CeilTabConsB & \CeilTabConsC & \CeilTabConsD & \CeilTabConsE & \CeilTabConsF & \GstarConsOne\\
High compute & \CeilCcHigh & \CeilFHigh & \CeilAminHigh & \CfullHigh & \CeilTauHigh & \MarginHigh & \CeilTabHighA & \CeilTabHighB & \CeilTabHighC & \CeilTabHighD & \CeilTabHighE & \CeilTabHighF & \GstarHighOne\\
\bottomrule
\end{tabular}

\smallskip
\begin{minipage}{0.95\linewidth}\footnotesize
Notes: ``Base'' is $(1-\tau)(1-c_c-f-a_{\min}-c_{\mathrm{full}})$, the ceiling at $G=0$, and $G^*$ is the market growth (\%) at which the ceiling reaches 100 with $k=0$. $c_c$, $f$ and $a_{\min}$ are scenario values, and $c_{\mathrm{full}}$ (points of revenue) is built in Table~\ref{tab:fullcost}. The high-compute case is within the value test, $c_c\le\CcValueTestCap\%$. Without the full-cost floors ($c_{\mathrm{full}}=0$), the base would be \MarginAccFav, \MarginAccCen, \MarginAccCons{} and \MarginAccHigh. The table's $\psi=0$ fits working capital and is a lower bound; for domestic development and equipment, $\psi=1$ under current law (\autoref{rem:taxes}) adds back $\tau kG$, at most \PsiAddMax{} points in the table, and development abroad adds back $\psi_{15}\tau kG$ with $\psi_{15}=\PsiFifteen$ in present value (in period cash on a steady growth path, $\psi_{15}$ evaluated at $G$, \PsiFifteenCashThirty{} at 30\%). Replacing $k=0.5$ at $\delta_K=0.25$ would subtract a further $(1-\psi\tau)\delta_Kk$ (\autoref{cor:ceiling}): $\delta_Kk=12.5$ points at the table's $\psi=0$, and \ReplCen{} points in the central case if the replacement equipment takes the 100\% allowance ($\psi=1$). Values computed by \texttt{code/calibrate\_horizons.py} and checked by \texttt{code/verify\_theory.py}.
\end{minipage}
\end{table}

\emph{The ceiling.} Table~\ref{tab:ceiling} evaluates \eqref{eq:Rmax}, holding output prices fixed. If compute absorbs \CcCen\% of revenue at full automation, a scenario value, the central margin term is \MarginCen{} points. An asset-light firm's ceiling is then \CeilCenFive, \CeilCenFifteen{} and \CeilCenThirty{} points for market growth of 5, 15 and 30\%, and $R_{\max}$ reaches 100 only if the market grows faster than about \GstarCenOne\% a year at these central inputs (\StGstarCcSixMin{} to \StGstarCcSixMax\% when the other costs move jointly over their stated ranges with compute held at \CcCen\%; Section~\ref{supp:stress} of the Supplementary Appendix). That threshold is $\tau+(1-\tau)(c_c+f+a_{\min}+c_{\mathrm{full}})$: \GdecTax{} points from taxes, \GdecCompute{} from compute, \GdecFees{} from fees, \GdecAcc{} from the accountability band and \GdecFull{} from the full-cost floors (maintenance \GdecMnt, security \GdecSec, AI failure \GdecErr{} and other \GdecAnc). Favorable parameters raise the ceiling to \CeilFavFive, \CeilFavFifteen{} and \CeilFavThirty{} and lower the threshold to about \GstarFav\%; conservative ones give \CeilConsFive, \CeilConsFifteen{} and \CeilConsThirty{} and \GstarCons\%; and in the high-compute case the margin term falls to \MarginHighR{} points and the threshold rises to \GstarHigh\%. Counting compute, fees and the accountability floor but none of the full-cost floors of running AI (\autoref{def:fullcost}), the four thresholds would be \GstarAccFav, \GstarAccCen, \GstarAccCons{} and \GstarAccHigh\%. The sensitivities of the central threshold rank the parameters:
\begin{itemize}
\item \emph{Compute share.} Over its scenario range of \CcLo{} to \CcHi\% of revenue, the threshold runs from \GstarCcLo{} to \GstarCcHi\% (\GstarAccCcLo{} to \GstarAccCcHi\% without the full-cost floors). This is the dominant uncertainty.
\item \emph{AI failure.} From its favorable to its conservative value, the failure term moves the threshold from \GstarCerrFav{} to \GstarCerrCons\%; with today's highest measured error rates and losses equal to the in-house cost of the work ($c_{\mathrm{err}}=\CerrHigh$ points), it would reach \GstarCerrHigh\% with no review, and \GstarCerrHighAi\% ($c_{\mathrm{err}}=\CerrHighAi$) with AI review at its bound.
\item \emph{Taxes and the accountability floor.} Over tax rates of 15 to 21\% the threshold runs from \GstarTauLo{} to \GstarTauHi\%, and at a 9\% effective rate it would be \GstarTauLow\% (\GstarAccTauLow\% without the full-cost floors). Over the band of $a_{\min}$ it runs from \GstarAminLo{} to \GstarAminHi\%.
\item \emph{Capital.} A capital intensity of 1.5 lowers the ceiling at 30\% market growth by \KcapLossNone{} points if the capital is not deductible, by \KcapLossDomLo{} to \KcapLossDomHi{} points if it is deducted when incurred, as domestic software development and equipment that qualifies for the 100\% allowance now are, and by \KcapLossForLo{} to \KcapLossForHi{} points in present value for development abroad (\autoref{rem:taxes}; about \KcapLossForCash{} points of period cash at 30\% growth, because deductions lag capital that grows faster than the discount rate). Owning the compute behind $c_c$ instead of buying it raises the cash ceiling only if the vendor's markup exceeds 3.3 times market growth when servers take the 100\% allowance of current law (\autoref{rem:owncompute}), about \OwnMarkupFive\% at 5\% market growth and \OwnMarkupFifteen\% at 15\%; a firm that elects out of the allowance, or whose servers are used mainly abroad, depreciates them over several years and faces a higher bar, between those values and the lower bound at $\psi=0$, for servers that cannot be deducted at all, of \OwnMarkupFiveNone\% and \OwnMarkupFifteenNone\%. The bar assumes owned capacity is fully used; idle capacity raises it.
\item \emph{Deferred revenue and a typical firm.} A firm that bills in advance has a negative cash capital intensity, which lowers the threshold: at $k=-0.3$ the central threshold is \GstarKdefCen\% (\GstarKdefFav\% and \GstarKdefCons\% in the favorable and conservative cases), and at the median ratio of the change in deferred revenue to the change in revenue among growing firm-years of 2023--25, \KdefMed{} (from SEC frames, which keep the latest filed value), it is \GstarKdefMedCen\%. A typical firm with \$1 to 5 billion of revenue in 2023--25 has a margin intercept \TypGapNet{} points below the central full-cost margin term and a threshold of \TypGstarK\% at its own capital intensity (Section~\ref{app:data-typical} of the Supplementary Appendix).
\end{itemize}

\emph{By horizon and with pass-through.} Along the scenario path of the compute share, the share moves with two prices. The price of a fixed level of performance falls by 5 to 10 times a year, while the cost of running frontier models on the hardest tasks has risen by 3 to 18 times a year, through token volume rather than price per token \citep{gundlach2025price}. The path prices the work that only frontier models reach near each horizon at a fraction of the labor cost it replaces and the rest at the falling price of fixed performance, so the central compute share rises to \CcCenTwentyEight\% in 2028 and then falls to \CcCenThirtyOne\% in 2031 and \CcCenThirtySix\% in 2036, with the favorable and conservative cases at half and twice these values (Section~\ref{supp:model-compute} of the Supplementary Appendix gives its two drivers). The headline compute share of \CcCen\% at 2025 prices is not a date on this path, which reaches it only in about \CcYearSix. The central threshold is then about \GstarCenTwentyEight\% in 2028, \GstarCenThirtyOne\% in 2031 and \GstarCenThirtySix\% in 2036, against \GstarCenOne\% at 2025 prices. These thresholds hold output prices fixed. If competition passes part of the fall in unit costs to customers, they rise. We let the industry pass the share $\varphi$ of every unit-cost saving to its price, the pass-through rate that the evidence measures, and evaluate the margin at the lower price. The saving is the top-tail anchor firm's own, from a unit cost of \PtCzero\% of price; a typical \$1--5 billion firm (\PtCzeroTyp\%) saves more, with thresholds of \PtGstarTypHalf{} and \PtGstarTypSeven\% (Table~\ref{tab:passthrough-recon}). All savings are exposed, the accountability floor and the costs of running AI included, because competition prices the full unit cost (\autoref{prop:entry}(d)). With $\varphi=0.5$ the price falls by \PtPriceCutHalf\% and the central threshold at 2025 prices rises to \GstarPhiHalfOne\%; with $\varphi=0.7$ it falls by \PtPriceCutSeven\% and the threshold is \GstarPhiSevenOne\%. Exposing only the saving in cost of revenue gives \PtGstarVarHalf{} and \PtGstarVarSeven\% (Section~\ref{supp:model-compute} gives the thresholds by horizon, and Table~\ref{tab:passthrough-recon} reconciles the readings). The scenario values of $\varphi$ span the evidence: close to zero in the short run for differentiated software \citep{choi2025dma} and 0.7 for an industry-wide cost shock \citep{ganapati2020energy}. Applying that medium-run, marginal-cost rate to every unit cost rests on the long-run argument of \autoref{prop:entry}(d).

Three conclusions follow. First, at the full-automation cost structure, complete automation does not carry a mature firm past a cash-basis Rule of 100 unless its capital needs are negligible and its market grows faster than about \GstarCenOne\% a year, if compute absorbs \CcCen\% of revenue at 2025 prices and the other costs are at their central values; a firm whose costs are leaner than that structure can score more. Every threshold is conditional on the scenario compute share, the range over its scenario bounds is wide, and pass-through raises every threshold except near the top of the compute share's scenario range, where compute costs nearly as much as the labor it replaces and the full-automation unit cost exceeds today's (Section~\ref{supp:stress}). Second, at any market growth a software market has sustained, a Rule of 1000 is not a steady state. It requires roughly \PullMultLo{} to \PullMultHi{} times the demand pull of the best public software firms today, and \autoref{prop:budget} and \autoref{cor:transient} bound how long such a surplus can last: every point of score above the saturated level is paid for out of a finite budget of share headroom. Section~\ref{supp:rule1000} of the Supplementary Appendix works through the numbers. Third, under free entry into a market for a homogeneous product the margin intercept is not set by automation. With free entry and replication cost approaching zero, \autoref{prop:entry} gives $R=G+(r-G)(k+k_a)+\rho$ for the Rule of 40, where $k_a$ is customer capital per dollar of revenue. With market growth of 5\%, a cost of capital of 10\%, $k=0.5$ and no customer capital, a firm without a scarce input scores \EntryScore; customer capital adds a normal return on what it cost to acquire, not a rent. A rent cannot lift the score above the full-cost ceiling (\autoref{prop:entry}(d)): at the central cost structure of Table~\ref{tab:ceiling}, a firm with 5\% market growth and $k=0.5$ scores at most the central margin term plus 2.5 to 3 points, so a cash-basis score of 100 is out of reach whatever the rent. (This applies \autoref{prop:entry}(d) after tax: with profits and rents taxed at $\tau$, growth capital deducted at $\psi$ and customer acquisition, which \autoref{ass:entry} expenses when incurred, deducted in full, its bound is the after-tax $R_{\max}$ of \autoref{cor:ceiling} with capital $k$, less $(1-\tau)k_aG$; here $k_a=0$.) In a competitive market for a homogeneous product, a durable high score requires a rent on an input rivals cannot replicate, up to that ceiling, and the score does not distinguish such a rent from market power.

\subsection{A central path to 2036}\label{sec:horizons-ranges}
Table~\ref{tab:calib} traces the maximal score of a top-tail asset-light firm, one at the 95th percentile of the 2023--25 cash score, along a central path, and Table~\ref{tab:tornado} moves one input at a time. The path is a scenario, not a forecast, and it describes one firm type: the attainable frontier is taken at a fixed share (Section~\ref{sec:model-attainable}). The calibration is on an after-tax cash basis, so $m_0$ is the after-tax margin intercept of \autoref{rem:taxes}.

\emph{Anchor.} The data identify neither demand pull nor the spend elasticity, and a single intercept across firms is not usable either: an intercept that fits typical firms lies below the observed margins of many of the best ones. We therefore invert the model firm by firm for the top tail of observed scores, under the calibration rule of Section~\ref{sec:empirics-anchor}: the \AnchorN{} firm-years of \AnchorFirms{} firms at or above the 95th percentile of the 2023--25 cash score (\AnchorPthree{} points), without the \AnchorAcqExcl{} acquisition-affected firm-years of the tail, with market growth set at the sample's median growth of \AnchorG\%. On the after-tax cash basis $m=m_0-(1-\tau)\sigma-kg$, so each firm's intercept is $m_0=m+(1-\tau)\sigma+kg$, with $\tau=0.18$ and measured capital spending for $kg$, treated as not deductible ($\psi=0$); with the deductions of current law each intercept would be lower by $\tau kg$, at most \AnchorTaxGap{} points at the quartiles. Growth spend $\sigma$ is not observed. The central basis counts as growth spend only selling and marketing expense above a maintenance level, $\sigma=\max(s_{\rm SM}-b_M,0)$, where $s_{\rm SM}$ is the selling-and-marketing share and $b_M=\MaintBase\%$ is the median share, in the previous year, of firm-years growing at or below the year's median: the netting behind the spend ratio of \AlphaNet{} in Section~\ref{sec:empirics}. Of the \AnchorN{} anchor firm-years, \AnchorSmAboveBase{} spend above that level. Zero and the whole selling-and-marketing share bound $\sigma$. We split each score into a durable part $R_\infty=m_0+(1-k)G$ and a transient growth surplus $\Pi=R-R_\infty$ on each basis.

On the central basis the median anchor firm-year had a margin intercept of \AnchorMzeroNetMed{} points (interquartile range \AnchorMzeroNetLo{} to \AnchorMzeroNetHi; 90\% firm-cluster bootstrap interval of the median \AnchorMzeroNetCiLo{} to \AnchorMzeroNetCiHi) and a growth surplus of \AnchorPiNetMed{} points. If all selling and marketing expense buys growth, the median intercept is \AnchorMzeroMed{} points (interquartile range \AnchorMzeroLo{} to \AnchorMzeroHi) and the surplus \AnchorPiMed{} points, with a 90\% firm-cluster bootstrap interval of \AncPiUbEoneLo{} to \AncPiUbEoneHi; \AnchorPiUbPos{} of the \AnchorN{} firm-years scored above their durable level. A negative surplus is inconsistent with the model at $X=1$, because the surplus is a maximum over spend that includes zero spend; it means either that the firm acts on a weight on growth above one or that its selling-and-marketing share overstates its growth spend. If none of it buys growth, the median intercept is \AnchorMzeroLbMed{} points and the surplus \AnchorPiLbMed{} (\AncPiLbEoneLo{} to \AncPiLbEoneHi). Only the last statement is a conclusion: if selling spend does not buy growth, the best scores of 2023--25 exceed their durable level, with an interval above zero in \AncPiLbCount{} views (Section~\ref{sec:empirics-anchor}). How much of those scores was transient the data do not determine. The central median intercept is \AnchorGapCenNet{} points below the central full-cost margin term at 2025 prices (\MarginCenR), and even the higher, upper-bound intercept is \AnchorGapCen{} points below it and \AnchorGapFav{} points below the favorable one (\MarginFavR), so automation still has room to raise the durable scores of the best firms, but by tens of points, not hundreds.

\begin{table}[t]\centering\small
\caption{Central path of the maximal Rule-of-40 score of a top-tail asset-light firm, at full cost and attainable pace (free-cash-flow-after-SBC basis, $X=1$, output prices fixed), with 90\% anchor-sampling intervals}\label{tab:calib}
\resizebox{\textwidth}{!}{\begin{tabular}{lrrrrrrr}\toprule
 & \multicolumn{2}{c}{Program} & \multicolumn{2}{c}{Maintenance-netted (central)} & All S\&M buys growth & None buys growth & \\\cmidrule(lr){2-3}\cmidrule(lr){4-5}
Horizon & Realized $\bar b\,\mathcal E$ & Target $m_{0F}$ & $m_0$ & $R^*=R_\infty$ & $R_\infty$ (bound) & $R_\infty$ (bound) & Dip\\\midrule
2025 (anchor) & -- & -- & 31 [24, 35] & 43 & 55 & 37 & -- \\
2028 & 7 & 62 & 33 & 46 [39, 50] & 56 [49, 66] & 40 [34, 47] & 3.6 \\
2031 & 20 & 66 & 38 & 50 [44, 54] & 59 [53, 68] & 45 [40, 52] & 3.9 \\
2036 & 41 & 68 & 46 & 59 [54, 62] & 65 [61, 72] & 55 [51, 60] & 4.2 \\
\addlinespace\multicolumn{8}{l}{\emph{Anchor sample: 2025 medians and the 2036 central score}}\\
 & Firm-years (firms) & $m_0$, netted & $m_0$, all & $\Pi$, all buys growth & $\Pi$, none buys growth & 2036, netted & 2036, all / none\\\midrule
E1, calibration rule (central) & 13 (7) & 31 & 42 & $-$1 [$-$18.2, $+$16.4] & 20 [$+$8.1, $+$26.5] & 59 & 65 / 55 \\
E1, every firm-year of the tail & 22 (14) & 26 & 42 & $-$1 [$-$10.7, $+$15.7] & 22 [$+$12.9, $+$27.0] & 56 & 65 / 53 \\
E2, organic & 19 (10) & 32 & 45 & $-$7 [$-$17.0, $+$15.4] & 18 [$+$6.8, $+$26.0] & 59 & 67 / 56 \\
E2 plus imputed, $M=4$ & 22 (14) & 27 & 42 & 0 [$-$13.6, $+$14.5] & 20 [$+$11.3, $+$26.9] & 56 & 65 / 55 \\
E2 plus imputed, $M=8$ & 22 (13) & 27 & 42 & $-$1 [$-$13.7, $+$15.9] & 20 [$+$12.2, $+$27.1] & 56 & 65 / 55 \\
\bottomrule\end{tabular}}
\par\vspace{2pt}\footnotesize\raggedright Notes: Percentage points, rounded half up; brackets: 90\% firm-cluster bootstrap intervals from the anchor's sampling error (999 draws of the anchor block, intercept, market growth and maintenance base redrawn together; the program and the target are held fixed). The 2025 row inverts the model for the median of the 13 firm-years of the audited sample at or above the 95th percentile of the 2023--25 cash score, without acquisition-affected firm-years (the calibration rule; Section~\ref{sec:empirics-anchor}): on the after-tax cash basis $m_0=m+(1-\tau)\sigma+kg$ with $\tau=0.18$, capital treated as not deductible ($\psi=0$). Growth spend $\sigma$ is the selling-and-marketing share in excess of the maintenance level of 25.1\% of revenue (the median share of firm-years growing at or below the year median, Section~\ref{sec:empirics}; central), the whole share (all buys growth) or zero (none does); market growth at the anchor's median growth of 12.7\%. Later rows: the central path, a scenario, not a forecast. The intercept moves toward the central full-cost margin term $m_{0F}$ at each horizon's compute share, $m_0=m_{0A}+\bar b\,\mathcal E\,(m_{0F}-m_{0A})$, with $\bar b\,\mathcal E$ the share of the program's benefits realized by the horizon under central execution (the upper bound of \autoref{prop:expected}(e)); $R_\infty=m_0+(1-k)G$ (\autoref{prop:optimum}) with $k=0$ and $G$ held at the anchor's median growth. The path reports the durable score on every basis; the transient surplus enters Table~\ref{tab:tornado}. Dip: the cost of the transitions in flight for a firm still transforming, central basis. Lower panel: medians of the intercept and of the surplus $\Pi=R-R_\infty$ with 90\% firm-cluster bootstrap intervals, and the 2036 central path from each sample's median intercept.
\end{table}

\emph{Pace.} The intercept does not follow an assumed path. It moves from the anchor toward the full-cost margin term at the pace that \autoref{prop:expected}(e), the version for a divisible program, allows: $m_0(t)=m_{0A}+\bar b\,\mathcal E(t)(m_{0F}-m_{0A})$, where $\mathcal E(t)\le\min\{1,p^{\rm emb}\rho T_{\mathrm{par}}n(t_R)\}$ is the expected share of the automation program embedded by date $t$, $\rho$ the binding capacity as a share of the program a year, and $\bar b$ the share of benefits that embedded processes realize. The calibration uses this upper bound, which is attained through 2036 under central execution (Section~\ref{supp:model-clock} of the Supplementary Appendix). The calibration assumes that every process is technically automatable from the program's start; it has no gate on AI capability. Execution has three parts.
\begin{itemize}
\item \emph{Capacity.} Central execution changes the work of \PaceCapCen\% of the firm's employees a year, within a range of \PaceCapSlow{} to \PaceCapFast\% that an employer survey of training needs to 2030 \citep{wef2025futurejobs}, single restructurings and exit rates loosely bracket. Fast execution changes \PaceCapFast\%. Capacity maps one for one onto the program (Section~\ref{supp:model-execution} of the Supplementary Appendix). For a program of \PaceNproc{} processes, the scenario values for decision bandwidth and security review allow a faster pace in both cases, so change capacity binds (\autoref{prop:attainable}(c)).
\item \emph{Risk.} A share $\varpi=\PaceSuccCen$ of attempts creates some value in central execution and \PaceSuccFast{} in fast execution. These are not success rates: the central value rounds down the share of transformations in a survey that created value, which includes those that missed their targets \citep{bcg2020flipping}. Transitions embed \PaceEmbCen{} years and \PaceEmbFast{} year after they start, successes realize \PaceRealCen\% and \PaceRealFast\% of their targeted savings, and reorganizations freeze the program about \PaceFrzCen\% and \PaceFrzFast\% of the time.
\item \emph{Embedding.} These shares are outcome shares that already count transitions lost to leadership changes and changes of strategy, so the calibration reads the embedding probability of \autoref{prop:expected} directly from them, $p^{\rm emb}=\varpi$ (\PaceQCen{} and \PaceQFast), and does not apply the disruption hazard a second time ($\nu_D=0$ in part (e)). Each wave attempts as much of the program as change capacity allows, so the program embeds \PaceRateCen\% of itself a year in expectation under central execution and \PaceRateFast\% under fast execution.
\end{itemize}
Under central execution a top-tail firm realizes \PaceCenEight\% of the benefits of its program by 2028, \PaceCenOne\% by 2031 and \PaceCenSix\% by 2036, having embedded about \CenEmbEight, \CenEmbOne{} and \CenEmbSix\% of it. Fast execution realizes \PaceFastEight, \PaceFastOne{} and \PaceFastSix\% and its upper bound reaches the whole program in \PaceFastDone{} (the exact share reaches \PayDoneShare\% in the same year and approaches one geometrically after that). Only embedded transitions realize value, so the yield per attempt is $\varpi\bar b=\YieldModel$ of the targeted savings, close to the unconditional realized share of \YieldMck{} that two surveys imply. Read like for like, the realized share by 2036 is \ShareExec\%, and if only transformations that met their targets count, \ShareBcgWin\% (Section~\ref{supp:model-execution} of the Supplementary Appendix, which also decomposes the change from the first version).

A firm still transforming also carries the cost of the transitions in flight. If each costs its savings over the time to payback, \PaceTpb{} years under central execution (the embedding time of Table~\ref{tab:param-sources}) and one year under fast execution, the cash score is lower by about \CenDipEight{} to \CenDipSix{} points under central execution and by \PaceDipFast{} points in 2028 under fast execution while the program runs; at two years of savings the central dip is \DipTpbTwoEight{} to \DipTpbTwoSix{} points. No source measures the cost of a transition in years of its savings. Under central execution realized gains first exceed the dip after about \PayCrossCen{} years (\CrossTpbTwo{} at two years of savings). That flow crossing is not payback: the cumulative gains repay the cumulative dip only after \PayBackCen{} years, or \PayBackDiscCen{} discounted at \acqRwAnnual, and the cumulative effect is still negative at three and at five years (\PayHoldThreeCen{} and \PayHoldFiveCen{} points of revenue times years). Both statements are for central execution: with the other inputs moved jointly over their stated ranges the payback runs from \StPayCenMin{} to \StPayCenMax{} years, and the cumulative effect at five years turns marginally positive at a few points (Section~\ref{supp:stress} of the Supplementary Appendix). In present value the program's gains, which persist after its upper bound reaches the whole program in \PayDoneCen, are worth \PayRatioCen{} times the dip (\PayRatioDoneCen{} counting only gains until then), and the net present value is \PayNpvCen\% of a year's revenue, with revenue held constant. Fast execution pays back after \PayBackFast{} years (\PayBackDiscFast{} discounted). The program therefore fails a five-year cash-payback test under central execution, though not under fast execution. An owner who sells within five years still gains from it if the sale price reflects the benefits still to come, which on the central path are worth \PayExitFiveCen\% of a year's revenue at year five; it would decline the program only if a sale recovered less than \PayExitBreakCen\% of that value, or if it must repay the program's costs in cash within its horizon. Under slow execution the net present value is \PayNpvSlow\% of a year's revenue, so the program is not worth funding at any horizon. The survey rates, the capacity range and the turnover and reorganization rates are sourced or grounded by precedent; the program size, embedding times, realized shares, the rates at which new leaders keep a program and strategies change, and the transition costs are scenario values. The path counts only the automation program. For comparison, between fiscal 2022 and 2024 the median core firm with at least \$100 million of revenue raised its GAAP operating margin by \RevealedOpMed{} points (interquartile range \RevealedOpLo{} to \RevealedOpHi, $n=\RevealedN$) and its cash margin after SBC by \RevealedFcfMed{} points, through cost programs of every kind; by 2028 the central path raises the intercept by \CenRiseEight{} points.

\emph{Central path.} Table~\ref{tab:calib} combines the anchor, the full-cost target at each horizon's compute share and central execution, with output prices fixed, $k=0$ and market growth held at \CenG\%. The path reports the durable score on every basis. A transient surplus is scaled with demand pull only under the upper bound on growth spend, where the median anchor firm-year has none; the central basis gives the median firm-year a surplus of \AnchorPiNetMed{} points, which the path leaves out and Table~\ref{tab:tornado} treats as a sensitivity. On the central basis the maximal durable score at fixed prices is \CenREight{} in 2028 [\CenREightLo, \CenREightHi], \CenROne{} in 2031 [\CenROneLo, \CenROneHi] and \CenRSix{} in 2036 [\CenRSixLo, \CenRSixHi], as the intercept rises from \AnchorMzeroNetMed{} to \CenMzeroEight, \CenMzeroOne{} and \CenMzeroSix{} points; net of the dip of transitions in flight it is \CenREightNet, \CenROneNet{} and \CenRSixNet. The brackets are 90\% intervals from resampling the anchor firms alone, with the program and the target held fixed; with every scenario input moved jointly, 90\% of the draws of the 2036 score lie between \StNetSixPFive{} and \StNetSixPNinetyFive{} (Section~\ref{supp:stress} of the Supplementary Appendix). If all selling and marketing expense buys growth, the intercept starts \AnchorShiftNet{} points higher and the score is \CenREightAllSM, \CenROneAllSM{} and \CenRSixAllSM; if none of it does, the score is \CenRLbEight, \CenRLbOne{} and \CenRLbSix. Even a completed program would bring the durable score only to the target plus market growth, about \CenTargetSix{} plus \CenG{} points in 2036. With pass-through of $\varphi=0.5$ the central path is \CenRPhiHalfEight, \CenRPhiHalfOne{} and \CenRPhiHalfSix, and with $\varphi=0.7$ it is \CenRPhiSevenEight, \CenRPhiSevenOne{} and \CenRPhiSevenSix: pass-through lowers only the target, and under central execution the firm has realized less than half of its program by 2036 (\StRealMin{} to \StRealMax\% over the stated ranges of execution).

\begin{table}[t]\centering\small
\caption{One-at-a-time sensitivity of the central maximal score (tornado), 2028 and 2036}\label{tab:tornado}
\begin{tabular}{@{}p{0.62\linewidth}rrr@{}}\toprule
Input moved to the ends of its range & Low & High & Swing\\\midrule
\multicolumn{4}{@{}l}{\emph{2036; central path 59}}\\
Execution: slow and fast & 45 & 75 & 30 \\
Transient surplus of the upper-quartile firm-year (all selling and marketing buys growth), pull and $\alpha$ at their ends & 72 & 91 & 19 \\
Market growth $G$ over the interquartile range of 2023--25 growth & 51 & 67 & 16 \\
Selling and marketing: none and all of it buys growth & 55 & 65 & 11 \\
Full-cost target: conservative and favorable case & 52 & 63 & 11 \\
Execution yield: BCG win zone only ($\varpi=0.30$, $\bar b=0.67$) & 50 & 59 & 9 \\
Capital $k=0.5$, not deductible & 52 & 59 & 6 \\
Anchor intercept: lower and upper quartile & 55 & 61 & 6 \\
Pass-through $\varphi=0.7$ of all savings (fixed prices at the central value) & 54 & 59 & 5 \\
A firm still transforming (in-flight dip) & 54 & 59 & 4 \\
Anchor sample: every firm-year of the tail, E2, $M=4$, $M=8$ & 56 & 59 & 3 \\
\addlinespace\multicolumn{4}{@{}l}{\emph{2028; central path 46}}\\
Transient surplus of the upper-quartile firm-year (all selling and marketing buys growth), pull and $\alpha$ at their ends & 73 & 116 & 43 \\
Selling and marketing: none and all of it buys growth & 40 & 56 & 17 \\
Market growth $G$ over the interquartile range of 2023--25 growth & 38 & 54 & 16 \\
Execution: slow and fast & 43 & 53 & 10 \\
Anchor intercept: lower and upper quartile & 40 & 49 & 9 \\
Capital $k=0.5$, not deductible & 39 & 46 & 6 \\
Anchor sample: every firm-year of the tail, E2, $M=4$, $M=8$ & 41 & 47 & 5 \\
A firm still transforming (in-flight dip) & 42 & 46 & 4 \\
Full-cost target: conservative and favorable case & 44 & 47 & 2 \\
Execution yield: BCG win zone only ($\varpi=0.30$, $\bar b=0.67$) & 44 & 46 & 1 \\
Pass-through $\varphi=0.7$ of all savings (fixed prices at the central value) & 45 & 46 & 1 \\
\bottomrule\end{tabular}
\par\vspace{2pt}\footnotesize\raggedright Notes: Points. Each row moves one input from the central path of Table~\ref{tab:calib} (maintenance-netted basis, $k=0$, central execution, the central full-cost target, fixed prices, $\alpha=0.35$) to the ends of its range, holding the others at their central values; rows sorted by swing. Ranges: anchor quartiles (netted basis); market growth $G$ over the interquartile range of revenue growth of the core firm-years of 2023--25 (5.4--21.4\%, own data); for the surplus row, the upper-quartile firm-year under the upper bound on growth spend with pull $\mu$ and $\alpha$ at their ends (2028: $\mu$ 1--2; 2031: $\mu$ 0.75--1.5; 2036: $\mu$ 0.5--1.25); slow and fast execution of Section~\ref{sec:horizons-ranges}; the win zone of the transformation survey alone ($\varpi=0.30$) with the realized share of successful transformations; the conservative and favorable full-cost cases; pass-through at a rate $\varphi=0.7$ of all unit-cost savings, margins at the post-cut price; $k=0.5$ of capital per dollar of new revenue, not deductible ($\psi=0$), against the central $k=0$; the two bounds on growth spend; the four alternative anchor samples. One-sided rows compare one alternative with the central path. The surplus row adds the upper-quartile firm-year's transient surplus, which the median firm-year does not have, so both of its ends lie above the central path. All are scenario values.
\end{table}

\emph{Sensitivity.} In 2036 the largest swing in Table~\ref{tab:tornado} comes from \TornFirstSix: the score is \TornPaceSixLo{} under slow and \TornPaceSixHi{} under fast execution, and \TornBcgWinSixLo{} if only the win zone of the survey creates value. The other large swings come from the surplus of an upper-quartile firm-year, scaled by demand pull and $\alpha$ at their ends (\TornSurplusSixLo{} to \TornSurplusSixHi), from market growth over the interquartile range of the growth of 2023--25, \GBandLo{} to \GBandHi\% (\TornGSixLo{} to \TornGSixHi; over 5 to 30\% it would be \TornGOldSixLo{} to \TornGOldSixHi), and from the attribution of selling and marketing expense (\TornAttributionSixLo{} if none and \TornAttributionSixHi{} if all of it buys growth). In 2028 the largest swing comes from \TornFirstEight, which can lift an upper-quartile firm-year to \TornSurplusEightHi{} if demand pull doubles and $\alpha$ is 0.45; the model scales the surplus with pull only under the upper bound on spend, so this is an order of magnitude for one firm type, not a path. No single input in the table lifts the 2036 central score to 100, and moving all the inputs together over their stated ranges does not either: in the joint-parameter stress test of Section~\ref{supp:stress} of the Supplementary Appendix the netted 2036 score runs from \StNetSixMin{} to \StNetSixMax. That test also shows which claims of this section hold only under their stated conditions, and the text states those conditions where each claim appears.

\emph{Ownership.} In the cases computed here, ownership changes the pace, not the target. A public firm that holds to a guided dip of one point of revenue a year realizes \VarDipOneShare\% of the program's benefits by 2036 instead of \PaceCenSix\%, and \VarDipTwoShare\% at two points. Faster decisions alone change nothing: a variant with shorter sponsor latency has no effect while change capacity, not decision bandwidth, binds. We do not compute an owner-led execution case; the companion paper \citep{murthy2026transformation} derives the dip budgets and decision speeds of different owners, and Section~\ref{supp:model-execution} of the Supplementary Appendix gives further cases. These paths are lower than company-reported figures would suggest for two reasons: they are on an annual cash basis after SBC, and the fastest-growing AI firms of 2025 and 2026 are private and therefore outside the anchor.

\subsection{Levers}
The model turns the question ``what does it take to approach the limit?'' into conditions on parameters, each of which can be checked against a firm's own data. Six levers move the limit:
\begin{enumerate}
\item \emph{A scarce, non-replicable input} that supports a rent $\rho$: proprietary data with a feedback loop, distribution, a license, or network effects. Under free entry, without it the score falls toward $G+(r-G)(k+k_a)$. A data advantage is durable only if the data are renewed or the market tips: accumulating data has increasing returns early but diminishing returns in the long run \citep{farboodi2021model}, while data-driven markets can tip to a single firm \citep{prufer2021competing}.
\item \emph{Demand pull} $\lambda h$ well above that of the typical firm, which shows up as growth that does not require rising selling-and-marketing intensity.
\item \emph{Capital that earns its keep}: either $k$ near zero, or capital whose operating return exceeds the condition in \autoref{cor:substitution}, so that replacing labor with computing capacity raises rather than lowers the cash-basis score.
\item \emph{Pricing that captures value} as inference costs fall and as customers' headcount, and with it their seat counts, shrinks, so that falling unit costs raise $\gamma$ rather than being passed through to customers.
\item \emph{A small accountability floor} $a_{\min}$ and small full-cost floors. Certification of financial statements, audit, model-risk validation (the firm's own: the U.S.\ banking agencies' model-risk guidance of April 2026 places generative and agentic AI outside its scope; \citealp{fed2026sr262,occ2026bulletin13}) and liability insurance remain human-accountable costs; \autoref{prop:amin} shows that part of the floor is a requirement on the diversity of the control design and on its IT general controls rather than a headcount. The costs of maintaining, securing and supervising AI behave in the same way.
\item \emph{High market growth} $G$, meaning a position where spending on labor is converting into spending on software.
\end{enumerate}
A further factor moves the speed of approach rather than the limit: the capacity to execute. Change capacity, decision bandwidth and security review bound how many transitions a firm can run, and failures, rollbacks, leadership changes and reorganizations waste that capacity (\autoref{prop:attainable} and \autoref{prop:expected}). For a firm that acts on a decision weight that is highest at small share (\autoref{prop:dynamic}, under its hypothesis (H)), \autoref{cor:priorities} orders the levers it controls by stage, although which comes first at a given position also depends on the weight the owner acts on and on how much selling spend buys growth (Section~\ref{sec:levers-position}): while it has share to gain, the productivity of its growth engine (lever 2) and its retention, the firm-level counterpart of market growth (lever 6), move its score most; once it is near saturation, overhead and the accountability floor do (lever 5), while retention still raises the durable score by $X-k$ a point (\autoref{prop:levers-position}(a)); and capital-heavy automation is cheapest for slow growers (lever 3). No lever is sufficient on its own. Because any surplus above the durable score is transient, a sustained cash-basis Rule of 100 requires the levers jointly to push the durable score $R_\infty=m_0+(1-k)G$ above 100.

The evidence does not say which levers bind today. Among the best performers a point of growth has come with well under a point less margin, and the top of the score distribution fell between 2016--19 and 2023--25 while fewer firms grew fast; but the best margins of firms growing 20\% or more may have fallen as well, and the data cannot separate the two. They therefore do not show whether the supply of growth opportunity (levers 2 and 6) or the cost of buying growth is the constraint that binds. The decline of the upper tail with size is weak and not robust, so the data say little about saturation (lever 2, through $h$), and the lower cash-basis scores of capital-intensive firms are an accounting effect of deducting capital spending, not evidence on lever 3.

\subsection{Levers by position on the growth--margin plane}\label{sec:levers-position}
% The lever-table macros (sections/tables/levers_position_macros.tex, generated by code/playbook/run_archetypes.py) are loaded in paper1/p1_preamble.tex.

The envelope results of Section~\ref{sec:model-statics} rank the effects of AI at a single point of the frontier. This subsection states the ranking by position and the one test of it that needs no estimate of the decision weight. It is not a decision rule for a particular firm: the companion paper \citep{murthy2026transformation} turns it into a playbook by position and organizational setup (ownership, readiness, regulated processes, labor law, product model, market structure and owner capability), adds research and development and cost of revenue as levers, and removes levers through gates that hold whatever the position.

Let the firm be at the Rule-of-$X$ optimum for its weight $X>k$, with $\bar\sigma=\infty$, spend $\sigma^*$, excess growth $\Delta^*$, growth $g^*=G_i+\Delta^*$ and growth surplus $\Pi$ (\autoref{prop:score}). Compare four proportional levers at the same proportion $x$: automating the share $x$ of the cost of the growth activities, which multiplies demand pull by $(1-x)^{-\alpha}$; raising demand pull by $x$; cutting a non-discretionary line of size $a$ (general and administrative cost, or maintenance R\&D) by $x$; and cutting human service cost $c_h$ by $x$.

\begin{proposition}[Levers by position]\label{prop:levers-position}
\begin{enumerate}
\item[(a)] To first order the four levers raise $R^*_X$ by $x\sigma^*$, $x\sigma^*/\alpha$, $xa$ and $x(1-\epsilon_p)c_h$. The last two also raise the saturated score $R_\infty$; the first two do not, and vanish as $s\to\bar s$. A point of retention $G_i$ raises $R^*_X$ and $R_\infty$ by $X-k$, and capital for labor raises $R^*_X$ if and only if $da>(g^*+\delta_K)\,dk$.
\item[(b)] Automating the growth activities outranks cutting a line of size $a$ if and only if
\[
X>X^\ddagger_0:=k+\frac{a^{1-\alpha}}{\alpha\ell},
\]
and raising demand pull outranks it if and only if $X>k+(\alpha a)^{1-\alpha}/(\alpha\ell)$. If the firm's observed spend $\hat\sigma$ is optimal at the weight $\hat X$, the first condition is $(X-k)/(\hat X-k)>(a/\hat\sigma)^{1-\alpha}$.
\item[(c)] The growth-engine gains are proportional to $h(s)^{1/(1-\alpha)}$, so a line of size $a$ overtakes them once headroom has fallen by the factor $(a/\sigma^*)^{1-\alpha}$.
\item[(d)] Cutting human service cost outranks a line of size $a$ if and only if $(1-\epsilon_p)c_h>a$; under free entry without a scarce input its long-run effect on the score is zero (\autoref{prop:entry}(b)).
\item[(e)] A plan that lowers $m_0$ by $D$ and raises $G_i$ by $\Delta G$ raises $R^*_X$ by exactly $(X-k)\Delta G-D$.
\item[(f)] Ranked at a screening weight $X^{\rm scr}$ instead of $X$, the growth-side levers of (a), (b) and (e) are understated and the others unchanged, so the ranking reverses whenever $X^{\rm scr}<X^\ddagger_0<X$.
\end{enumerate}
\end{proposition}

\noindent The proof is in Section~\ref{app:proofs} of the Supplementary Appendix.

\smallskip
\emph{What the ranking needs from the data.} The data do not identify the level of the weight $\hat X$ that rationalizes a firm's spend (Section~\ref{sec:empirics}), so the useful form of part (b) is its ratio test, which compares $(X-k)/(\hat X-k)$, the owner's net weight relative to the one its spending reveals, with $(a/\hat\sigma)^{1-\alpha}$ and needs only the observed overhead line, the observed growth spend and $\alpha$. At the revealed weight itself the growth engine outranks overhead exactly when the firm spends more on growth than on the overhead line. The switch weight $X^\ddagger_0=k+(\hat X-k)(a/\hat\sigma)^{1-\alpha}$, which also needs demand pull, its comparison with screening weights in part (f) and the retention rate $G_i=g-\hat\sigma/(\alpha(\hat X-k))$ depend on the level of $\hat X$, and we do not report them.

Table~\ref{tab:levers-position} evaluates the proposition at six positions of the 2023--25 core sample: the cell medians of non-acquisition firm-years in a growth corner ($g$ at or above its upper quartile, $m$ below its lower quartile), a balanced upper-quartile scorer, a mature cash generator (growth below the median, margin in the upper quartile), a sub-40 grower, a stalled firm (both below their medians) and the near-frontier elite of Section~\ref{sec:horizons} (score at or above the 95th percentile). Growth spend $\hat\sigma$ is, on the central basis of Section~\ref{sec:horizons-ranges}, selling and marketing above the maintenance level of \lpMaintBase\% of revenue, with all and none of selling and marketing as bounds, and $\alpha=0.3$, the illustrative value of Section~\ref{sec:model} (the central path of Section~\ref{sec:horizons-ranges} uses \CenAlpha); in the table's last column $\delta$ is the replacement rate $\delta_K$. The table shows only quantities that need no weight: the cell medians, the ratio $(a/\hat\sigma)^{1-\alpha}$ and the lever that ranks first at the position's own weight. No weight is estimated.

\IfFileExists{sections/tables/tab_levers_position.tex}{\begin{table}[t]
\centering
\caption{Levers at six positions of the 2023--25 core sample: the ratio test of Proposition~\ref{prop:levers-position}(b), which needs no growth weight}\label{tab:levers-position}
\scriptsize
\setlength{\tabcolsep}{3pt}
\begin{tabular}{@{}lrrrrrrrrllr@{}}
\toprule
 & & & & & & & & & \multicolumn{2}{c}{First at own weight} & \\
\cmidrule(lr){10-11}
Position & $n$ & $g$ & $m$ & S\&M & $\hat\sigma$ & G\&A & $k$ & $(a/\hat\sigma)^{1-\alpha}$ & central & all S\&M & $g+\delta$\\
\midrule
Growth corner & 28 & 29 & \ensuremath{-}23 & 50 & 25 & 24 & 0.09 & 0.97 & growth engine & growth engine & 0.54 \\
Balanced upper quartile & 43 & 21 & 14 & 26 & 1 & 10 & 0.10 & 3.94 & G\&A cut & growth engine & 0.46 \\
Mature cash generator & 52 & 8 & 20 & 19 & 0 & 10 & 0.19 & -- & G\&A cut & growth engine & 0.33 \\
Sub-40 grower & 48 & 17 & \ensuremath{-}5 & 43 & 18 & 19 & 0.12 & 1.05 & G\&A cut & growth engine & 0.42 \\
Stalled & 84 & 2 & \ensuremath{-}12 & 31 & 6 & 20 & 1.21 & 2.23 & G\&A cut & growth engine & 0.27 \\
Near-frontier elite & 13 & 36 & 22 & 20 & 0 & 15 & 0.04 & -- & G\&A cut & growth engine & 0.61 \\
\bottomrule
\end{tabular}

\smallskip
\begin{minipage}{0.98\linewidth}\scriptsize
Notes: cell medians of non-acquisition core firm-years, fiscal 2023--25, in percent of revenue except $n$, $k$ and the ratio; $m$ is the free-cash-flow margin after stock-based compensation; $\alpha=0.3$. $\hat\sigma$ is growth spend on the central maintenance-netted basis: the selling-and-marketing share less a maintenance base of 25.1\% of revenue, the median S\&M share, in the previous year, of firm-years growing at or below the year's median, and zero below it (the rest of S\&M is a cost line); $a$ is the general and administrative share. $k$ is capital spending and capitalized software over revenue divided by $g$. The ratio $(a/\hat\sigma)^{1-\alpha}$ equals $(X^\ddagger_0-k)/(X-k)$ at the weight $X$ at which the position's spending is optimal: a 10\% cut in general and administrative cost outranks a 10\% saving in the cost of growth at every weight whose distance from $k$ is below the ratio times that $X-k$, and at the position's own weight the cut ranks first if and only if the ratio exceeds one; with no growth spend (--) the cut ranks first at every weight. ``Central'' is the lever that ranks first at the position's own weight on that basis; ``all S\&M'' is the same test with all selling and marketing counted as growth spend, a bound under which the ratio lies between 0.51 and 0.81; with no S\&M counted as growth spend, the other bound, the cut ranks first at every position. No weight is estimated (Section~\ref{sec:empirics}). $g+\delta$ is the opex saved per dollar of compute capital needed to raise the score at $\delta=0.25$. Generated by \texttt{code/playbook/run\_archetypes.py}.
\end{minipage}
\end{table}
}{}

The ranking depends on how much selling spend buys growth. On the central basis and at the weight a position's own spending reveals, cutting a share of general and administrative cost outranks automating the same share of the growth engine at \lpNGACutFirstNet{} of the \lpNPositions{} positions of Table~\ref{tab:levers-position} at the central maintenance level; over the 90\% bootstrap interval of that level the growth engine leads at \StLeverMin{} to \StLeverMax{} positions (Section~\ref{supp:stress} of the Supplementary Appendix). The growth engine leads only in the \MakeLowercase{\lpGEFirstNetList}, whose selling and marketing is far above the maintenance level, and \lpNNoGrowthSpend{} positions (\MakeLowercase{\lpNoGrowthSpendList}) spend nothing above it. Where a position has growth spend, the ratio runs from \lpRatioMin{} to \lpRatioMax; by part (b) the growth engine ranks first for an owner whose net weight exceeds that ratio times the revealed one. If all selling and marketing buys growth, the growth engine leads at \lpNGEFirstAllSM{} of the \lpNPositions{} positions, with ratios from \lpRatioAllSMMin{} to \lpRatioAllSMMax; if none does, at \lpNGEFirstNoSM{} positions. This is a ranking against the growth engine, not of every lever: the companion paper, which also values research and development as a cost line, ranks engineering above both at all \lpNPositions{} positions on the central basis \citep{murthy2026transformation}. Whether a firm ranked at a screening weight should cut overhead first depends on the level of the revealed weight, which the data do not identify (Section~\ref{sec:empirics}): over the range of $X-k$ that the proxies imply, from about \XkLo{} to \XkHi, the switch weight $X^\ddagger_0$ can lie on either side of the screening weights, so the table does not answer that question. The gains from the growth engine are transient and those from overhead durable (part (a)), and capital for labor faces a hurdle that rises with growth (part (a)). Overhead here is general and administrative cost; the table does not rank research and development or cost of revenue.

